\lesson{5}{Mar 07 2022 Mon (13:56:18)}{Fascination with Fear}{Unit 3}

\subsubsection{The Fascination with Fear}
\label{sub_sub_sec:the_fascination_with_fear}

The \textbf{suspense}, \textbf{mystery}, and \textbf{terror} that some stories
produce have attracted readers for hundreds of years.
Here's a timeline:

\paragraph{1765}
\label{par:1765}

\textbf{Horace Walpole} writes \textit{The Castle of Otranto}, the first, and
extremely popular, Gothic novel. Gothic fiction's goal is to inspire terror and
explore the supernatural; it contains elements of horror and romance.

% paragraph 1765 (end)

\paragraph{1800}
\label{par:1800}

The first known \textbf{Vampire Romance}, \textit{"Wake Not the Dead,"} by
\textbf{Johann Ludwig Tieck}, is published.

% paragraph 1800 (end)

\paragraph{1818}
\label{par:1818}

\textbf{Mary Shelley's} novel \textit{Frankenstein:} or, \textit{The Modern
Prometheus} is the first monster story.

% paragraph 1818 (end)

\paragraph{1831}
\label{par:1831}

\textbf{Victor Hugo} publishes \textit{The Hunchback of Notre Dame}, considered
a truly Gothic novel, with its bells and gargoyles.

% paragraph 1831 (end)

\paragraph{1839 to 1846}
\label{par:1839_to_1846}

\textbf{Edgar Allan Poe} publishes short stories with all things Gothic: death,
decay, insanity, ghosts, and live burials.

% paragraph 1839_to_1846 (end)

\paragraph{1840}
\label{par:1840}

Influenced by Gothic fiction and the ideas of the time, a new literary period
emerges: dark romanticism. It is a reaction against the optimistic ideas of the
literary movement of transcendentalism. Dark romanticism questions humanity's
goodness and is fascinated by darkness, the supernatural, and social and moral
decay.

% paragraph 1840 (end)

\paragraph{1885}
\label{par:1885}

\textbf{Robert Louis Stevenson} publishes \textit{The Strange Case of Dr.
Jekyll} and \textit{Mr. Hyde}, the classic story showing the darker side of man.

% paragraph 1885 (end)

\paragraph{1897}
\label{par:1897}

\textbf{Abraham Stoker} publishes \textit{Dracula}.

% paragraph 18971897 (end)

As you can see, the $1800$ s saw the development of many plot elements that
continue today.

\begin{enumerate}
  \label{enum:plot_elements_that_continue_today}

  \item Vampires
  \item Monsters
  \item Ghosts
  \item Crazed men and women.
\end{enumerate}

The expansion of this fear-loving literary genre was due in part to the society
at the time, which feared and, yet, was fascinated by death, the supernatural,
and humanity's dark side.

Authors began to explore rational fears in society through the irrational fears
portrayed in horror fiction.

\begin{definition}[Rational Fear]
  \label{def:rational_fear}

  A \textbf{Rational Fear} is a fear of something that actually poses some
  danger, such as diseases, major life changes, and accidents.
\end{definition}

\begin{definition}[Irrational Fear]
  \label{def:irrational_fear}

  An \textbf{Irrational Fear} is a fear of something that cannot actually do
  harm, such as mice, monsters, etc.
\end{definition}

\begin{example}
  \label{exm:irrational_fear}

  Frankenstein and Dr. Jekyll and Mr. Hyde are prime examples of the common
  rational fear of scientific advances and experiments, both works of fiction
  show an experiment going wrong, with the creation of a monster that torments
  society.
\end{example}

% subsubsection the_fascination_with_fear (end)

\subsubsection{Themes that Make Us Scream}
\label{sub_sub_sec:themes_that_make_us_scream}

Some literary critics find stories that evoke the most basic human emotion of
fear to be sensational, superficial, and without literary merit. A survey of
authors reveals that \textbf{Gothic Horror} explores what it means to be human.
It communicates something about life, making the thrills and chills meaningful.

Scary stories have a purpose beyond entertainment. Their themes communicate
something about life and humanity.

\begin{note}
  \label{nte:scary_stories_theme}

  A theme is never one word; it should always be stated as a sentence.

  \begin{itemize}
    \label{item:themes}

    \item Difficult experiences make us stronger.
    \item People should not take advantage of one another.
    \item Love is a sweet and bitter experience.
    \item Fear can keep us from trying new things.
    \item Humanity desires peace but cannot help engaging in conflict.
    \item A friend is the most valuable possession one can have.
    \item Loneliness helps us discover who we are as individuals.
    \item Excessive pride can destroy a person.
  \end{itemize}
\end{note}

% subsubsection themes_that_make_us_scream (end)

\subsubsection{Techniques that Make Us Tremble}
\label{sub_sub_sec:techniques_that_make_us_tremble}

The following devices are also used by authors to create suspense.

\paragraph{Sentence Structure}
\label{par:sentence_structure}

\begin{itemize}
  \label{item:sentence_structure}

  \item \textbf{Rhythm}: A combination of long and short sentences that draw the
    reader into the experience of the story.
  \item \textbf{Pace}: The use of phrases, sentences, and punctuation used to
    build tensions and suspense.
  \item \textbf{Anaphora}: The repetition of a word or phrase at the beginning
    of two or more successive clauses to emphasize an idea or create suspense.
\end{itemize}

% paragraph sentence_structure (end)

\paragraph{Time}
\label{par:time}

\begin{itemize}
  \label{item:time}

  \item \textbf{Foreshadowing}: The use of hints or clues that suggest future
    events.
  \item \textbf{Flashback}: An interruption in the sequence of time where a past
    event is presented in the midst of the portrayal of a current event.
\end{itemize}

% paragraph time (end)

\paragraph{Irony}
\label{par:irony}

\begin{itemize}
  \label{item:irony}

  \item \textbf{Verbal}: A device in which one means does not match with what
    one says.
  \item \textbf{Situational}: A device in which what one expects to happen does
    not actually occur.
  \item \textbf{Dramatic}: A device in which the audience knows something that
    the characters don't.
\end{itemize}

% paragraph irony (end)

\paragraph{Point of View}
\label{par:point_of_view}

\begin{itemize}
  \label{item:point_of_view}

  \item \textbf{First Person Narrators} are characters in the story and
    participate in action.
  \item \textbf{Second Person Narrators} address the reader as if he or she were
    part of the story's action.
  \item \textbf{Third Person Omniscient Narrators} are outside the action but
    can give insight into the thoughts of all characters.
  \item \textbf{Third Person Limited Narrators} are outside the action of the
    story but relate the thoughts, feelings, and experiences of one or two
    characters.
  \item \textbf{Third Person Objective Narrators} are outside the story. They
    look on watching every detail like a fly on the wall, but they cannot tell
    us what a character is thinking. They can only tell us what they observe.
\end{itemize}

% paragraph point_of_view (end)

\begin{definition}[Taphophbia]
  \label{def:taphophbia}

  \textbf{Taphophbia} is the fear of being buried alive. At the time, Edgar
  Allan Poe was writing, however, live burials did occur and were a common fear
  among many. Inventions such as the safety coffin were even developed to help
  people overcome their fear of being buried alive. According to one English
  reformer, prior to the $1900$ s, $149$ actual live burials took place in the
  previous century.
\end{definition}

% subsubsection techniques_that_make_us_tremble (end)

\newpage
